\documentclass[11pt,twoside]{article}
\usepackage{fontspec}
\usepackage{literary-edition}
\usepackage{critical-edition}
\usepackage{polyglossia}
\usepackage[paperwidth=16.50cm,paperheight=23.00cm,top=2.20cm,inner=2.65cm,outer=1.90cm,bottom=2.20cm]{geometry}
\setmainlanguage{spanish}
\setmainfont{Paskerville}[RawFeature=lnum]
\usepackage[breaklinks=true,hidelinks]{hyperref}
\usepackage[center,info,cam,a4]{crop}
\crop[cam]
%
\title{\fontsize{40}{20}\selectfont\textbf{HOMÉRICA}}
\date{\vspace{40pt} \textsc{Traducción fragmentaria de las obras homéricas por}\\\vspace*{14pt}{\huge\textsc{\textbf{Leopoldo Lugones}}}\\\fontsize{60}{11}\selectfont{\vfill\textbf{\texttt{\LaTeX}}}}
%
\newcommand{\ehuno}{\emph{\textbf{E.-H.\textsuperscript{(1)}}}}
\newcommand{\ehdos}{\emph{\textbf{E.-H.\textsuperscript{(2)}}}}
\newcommand{\ehter}{\emph{\textbf{E.-H.\textsuperscript{(3)}}}}
\newcommand{\ehctr}{\emph{\textbf{E.-H.\textsuperscript{(4)}}}}
\newcommand{\neh}{\emph{\textbf{N. E.-H.}}}
%
\newcommand{\separator}{%
	\par\noindent
	\hbox to \linewidth{\hspace*{15pt}\cleaders\hbox{.\quad}\hfil\hspace*{40pt}}%
	\par
}
%
%
\usepackage{metalogo}
\renewcommand{\footTmark}[1]{\textsuperscript{*}}
\renewcommand{\footTcall}[1]{}
\linenummargin{outer}
\nummarginsep{-12pt}{-40pt}
\thepagelook{upter}
\setlength{\pwidth}{0.90\textwidth}
%
%
\begin{document}
\pagenumbering{gobble}
	\maketitle
\newpage
%
\null
\vfill
\begin{flushright}
	Edición y composición en \texttt{\XeLaTeX} por \textsc{Ioannes Bracmanus} (2025).
\end{flushright}
\newpage
%
\tableofcontents
\newpage
%
\pagenumbering{Roman}
\setcounter{page}{1}
\simplehead{\textsc{Prefacio}}
\section*{\centering\textsc{Prefacio}}
\addcontentsline{toc}{section}{\textsc{Prefacio}}\par
%
\noindent\textsc{Las traducciones}, en el área de la literatura, además de tener su valor filológico o científico por su precisión, tienen valor estético. El gran \textsc{Leopoldo Lugones}, probablemente el mayor escritor que hubo en la historia de la literatura argentina, tradujo parte de las obras homéricas, desde episodios enteros hasta apenas unas líneas; y tal es su valor estético que merecen ser editadas por lo mismo, como cualquier traducción. El traductor, además de eximio helenista, fue un un auténtico compositor, un genio del \emph{verbo poético}; por lo mismo, además de imponerse la tarea de traducir fielmente a Homero desde su lengua, trasladando las acepciones al castellano con la mayor fidelidad posible, se impuso también la tarea de trasladar el ritmo, la forma. Para la traducción Lugones adopta el metro alejandrino por considerarlo, en su extensión y estructura, el más cercano al hexámetro dactílico griego, y es así que traduce la primera línea como «Canta, diosa, el encono de Aquiles Peleyades», a partir del mítico «Μῆνιν ἄειδε θεὰ Πηληϊάδεω Ἀχιλῆος».\par
%
El sistema de siglado para agilizar las referencias es el siguiente: \ehuno → \emph{Estudios Helénicos I: La funesta Helena} (1923); \ehdos → \emph{Estudios Helénicos II: Un paladín de la Ilíada} (1923); \ehter → \emph{Estudios Helénicos III: La dama de la Odisea} (1924); \ehctr → \emph{Estudios Helénicos IV: Héctor el domador} (1924); \neh → \emph{Nuevos Estudios Helénicos} (1928). Todos fueron editados por la Biblioteca Argentina de Buenas Ediciones Literarias (Babel).\par
\newpage
%
\pagenumbering{gobble}
\section*{\raggedleft\fontsize{30}{14}\textsc{Ilíada}}
\addcontentsline{toc}{section}{\emph{\textsc{Ilíada}}}
\repage
%
\pagenumbering{arabic}
\setcounter{page}{3}
%Canto 1.
\include{Rapsodias/1}
\repage
%
%Canto 2.
\setmainhead{\textsc{\emph{Ilíada}, Canto II}}
\bnum
\section*{\centering\textbf{CANTO II}\footT{\neh, pp. 92-93.}}
\addcontentsline{toc}{subsection}{\textsc{Canto II}}
%
\separator
%
\postart
\setotherline{474}Tal como los cabreros distinguen fácilmente\\
las copiosas majadas, si en el pasto se mezclan,\\
así les dan los jefes el orden conveniente\\
para el combate. Entre ellos va el rey Agamenón,\\
parecido por ojos y testa al dios tonante;\\
por el talle a Ares y por el pecho a Poseidón.\\
Cual la mayor cabeza del hato es el pujante\\
toro que entre las vacas sobresale, aquel día\\
dispuso Zeus que Atrida fuera el más arrogante\\
de porte entre los muchos héroes que allí había.\par
\poend
\enum
%
\separator
\repage
%
%Canto 3.
\include{Rapsodias/3}
\repage
%-
%Canto 4.
\include{Rapsodias/4}
\repage
%
%Canto 5.
\include{Rapsodias/5}
\repage
%
%Canto 6.
\include{Rapsodias/6}
\repage
%
%Canto 7.
\setmainhead{\textsc{\emph{Ilíada}, Canto VII}}
\bnum
\section*{\centering\textbf{CANTO VII}\footT{\ehdos, p. 101.}}
\addcontentsline{toc}{subsection}{\textsc{Canto VII}}
%
\separator
\postart
\setotherline{400}
¡Nadie acepte el tesoro de Alejandro, ni a Helena!\\
Pues un niño inocente vería cuán cercanas\\
están de los troyanos, perdición, y condena.\par
\poend
\separator
\enum
\repage
%
%Canto 8.
\setmainhead{\emph{\textsc{Ilíada}}, Canto VIII}
\bnum
\section*{\centering CANTO VIII\footT{\ehdos, pp. 102-104.}}
\addcontentsline{toc}{subsection}{\textsc{Canto VIII}}
%
\separator
\postart
\setotherline{130}
Y allá entre mil estragos el desastre ocurriera,\\
y los de Ilión quedaran cual corderos cercados,\\
si al punto el padre de hombres y dioses no lo viera.\\
Tronando, pues, horrísono, lanzó a tierra el fulgente\\
rayo ante los corceles de Diomedes que, entonces,\\
al surgir la terrible llama de azufre ardiente,\\
recularon de espanto bajo el carro de bronce;\\
y escapándose a Néstor el rendaje luciente,\\
\dlog{con miedo en su alma díjole a Diomedes:}{---¡Tidida,}\\
Revolvamos los fuertes potros para escaparnos!\\
¿No ves que el triunfo niégate Dios, y que hoy Zeus Kronida\\
otorga a ése la gloria que otra vez querrá darnos?\\
Pues por bravo que sea, no habrá varón que acierte\\
a estorbar el designio de Dios, mucho más fuerte.\par
Diomedes el cumplido guerrero, así le dijo:\\
---Sí, anciano, lo que hablaste, justo y cuerdo es, de fijo;\\
mas, pecho y alma embárgame cruel dolor, porque un día,\\
dirigiendo Héctor a los troyanos la palabra,\\
dirá que de él Tideides hacia la flota huía.\\
¡Y antes que así se alabe, que la tierra se me abra!\par
Mas, respondióle Néstor Gerenio el caballero:\\
---¡Ay de mí, qué has dicho hijo de Tideo el guerrero!\\
Aunque Héctor vil y flojo te llame, tales cosas\\
ni troyanos ni dárdanos creerán, ni las esposas\\
troyanas cuyos bravos y equipados maridos,\\
tú, en la flor de los años, entre el polvo has tendido.\par
Así habló, y revolviendo los corceles gallardos,\\
entre el tumulto huyeron; mas, con ruido inquietante,\\
hacían los troyanos y Héctor llover sus dardos;\\
y alzó el grito el grande Héctor del casco tremolante:\par
---¡Tideides, si los dánaos de la caballería,\\
grande honra con el solio, mesa y brindis te hacían,\\
van a despreciarte ahora que como mujer corres!\\
huye, pues, vil muchacha! Ni escalar nuestras torres\\
te permitiré, ni hacia los navíos llevarte\\
las mujeres, pues previo castigo pienso darte!\par
Dijo, y dudó Tideides entre seguir la huída,\\
o tornar los caballos para pelear de frente.\\
Tres veces vacilaron su corazón y mente,\\
y otras tantas el próvido Zeus, desde el monte Ida,\\
tronó por los troyanos, en signo fehaciente\\
de la victoria al bélico azar indefinida.\par
\poend
\separator
%
\postart
\setotherline{191}
¡Seguidme al punto, y rápidos, veamos de arrancar\\
su escudo a Néstor, célebre en la propia mansión\\
celeste, porque es todo de oro, hasta la armazón;\\
y al jinete Diomedes de los hombros sacar\\
el rico arnés que Hefesto le cinceló; pues creo\\
que si logramos esas dos cosas, los aqueos\\
esta misma noche hácense en su flota a la mar!\par
\poend
\separator
%
\postart
\setotherline{253}
No hubo quien de los muchos dánaos aventajara\\
al Tideida, en lanzarse con los raudos corceles\\
fuera del foso, a objeto de pelear cara a cara.\par
\poend
\separator
%
\postart
\setotherline{492}
Y aquéllos desmontaron para escuchar entonces\\
lo que les decía Héctor caro a Zeus, empuñando\\
su lanza de once codos cuya punta de bronce\\
que un aro de oro afianza, se erige relumbrando.\\
Y así, apoyado en ella, los fué aquél arengando.\par
\poend
\separator
%
\postart
\setotherline{529}
Hagamos, pues, la guardia nocturna, y bien temprano\\
salgamos pertrechados de todas armas, y ante\\
la flota, al ardiente Ares convoquemos ufanos.\\
Y sabré si Diomedes, el Tideida pujante,\\
desde la escuadra al muro me arroja, o si lo mato\\
con el bronce, y sangrientos despojos le arrebato.\\
Mañana ha de probarnos su valor con su aguante\\
cuando lo ataque a lanza; mas, creo que adelante\\
caérá herido junto con muchos camaradas\\
al salir el sol.\par
\poend
\separator
\enum
%
%Canto 9.
\setmainhead{\emph{\textsc{Ilíada}}, Canto IX}
\bnum
\section*{\centering\textbf{CANTO IX}\footT{\ehdos, pp. 73-74.}}
\addcontentsline{toc}{subsection}{\textsc{Canto IX}}

\separator
\postart
\setotherline{29}
Dijo, y todos inmóviles y en silencio han quedado.\\
Largo tiempo afligida calló la gente aquea\\
hasta que al fin Diomedes el buen jefe así ha hablado:

---Atrida, rechazo ante todo tu absurda idea,\\
pues tal es mi derecho de rey en la asamblea.\\
Mas, no te enojes, aunque mi honra heriste primero,\\
dándome entre los dánaos por flojo y mal guerrero;\\
lo que, de todo argivo, mozo y viejo, es sabido.\\
Si el hijo del artero Kronos, entre dos cosas,\\
darte el honor supremo del cetro ha decidido,\\
te negó el valor que es la fuerza más poderosa.\\
¡Desdichado! ¿Creíste que es tan floja y de escaso\\
valor la gente aquea como dices? Si acaso\\
al regreso te incita tu ánimo, parte ya,\\
que abierta senda tienes y junto al mar está\\
la gran flota que desde Micenas te siguier.\\
Mas los otros aqueos de larga cabellera,\\
hasta que a Troya hundamos se quedarán. Y si\\
también al caro suelo patrio huyen navegando,\\
yo y Esténelo, solos, quedaremos peleando\\
Hasta el fin de Ilión, puesto que dios nos trajo aquí.\par
\poend
\separator
%
\postart
\setotherline{308}
Divino Laertiades, Ulises ingenioso\\
preciso es que os exponga con veraz desembozo\\
lo que yo pienso y cómo lo he de cumplir, y dónde,\\
para que ceséis vuestro circunloquio enfadoso.\\
Pues, igual que las puertas del Hades, me es odioso\\
quien dice lo contrario de lo que en su alma esconde.\par
\poend
\separator
%
\postart
\setotherline{340}
¿Acaso, entre los hombres, sólo aman los Atridas\\
a sus esposas? Todo noble y cuerdo varón,\\
quiere y cuida la suya, cual yo, de corazón\\
amé a la que tuve, aunque por las armas fué habida.\par
\poend
\separator
%
\postart
\setotherline{697}
Agamenón rey de hombres, gloriosísimo Atrida,\\
nunca al irreprochable Peliónida debiste\\
suplicar, prometiéndole donación tan crecida;\\
pues si ya era orgulloso, mucho más lo engreiste.\\
Por lo que hace a él, dejémoslo que se vaya o se quede.\\
Cuando le dé la gana, de nuevo peleará,\\
o cuando un dios lo incite; pero lo que procede,\\
yo os lo diré y lo haremos: id a acostaros ya,\\
después que hayan saciado vuestro apetito, ahora,\\
comida y vino, que esto fuerza y valor nos dá;\\
y al despuntar la bella rosodáctila Aurora,\\
forma tú ante las naves gente y carros ligeros,\\
y a todos animando, lucha entre los primeros.
\poend
\enum
%
%
\end{document}